\documentclass[11pt]{ctexart}
\usepackage[a4paper,margin=1in]{geometry}
\usepackage{graphicx}
\usepackage{xcolor}
\usepackage{hyperref}
\definecolor{refgray}{gray}{0.45}
\title{\textbf{市场最新观点汇总}\\[0.4em]{\large 全球市场主线从需求侧叙事转向供给侧结构性分化，储能、铜、稀土等板块的定价逻辑正在重塑。}}
\author{KC桌面 自动生成}
\date{260601}
\begin{document}
\maketitle
\section*{一页摘要}
\begin{itemize}
  \item 全球储能产业正经历供给端结构性重组，而非简单的需求回暖，企业从规模竞赛转向成本与场景的双重分化。
  \item 电池行业当前周期与2021-2022年超级周期本质不同，盈利可持续性而非价格弹性成为关键，市场焦点应从价格修复转向出货量驱动的现金流增长。
  \item 铜价核心支撑来自非美市场供给收紧与美国进口激增导致的库存地理错配，而非AI等需求故事，矿山复产延迟是未来两年关键变量。
  \item G10消费者健康度分化明显，北美实际收入增长放缓至历史低位，消费韧性正从收入驱动转向储蓄消耗，欧元区相对稳健。
  \item 中东局势缓和情景下，市场低估了南非等具备高油价负相关、低利率敏感与超卖三重特征市场的反弹弹性。
  \item 人民币中间价模型显示政策当局正系统性引入模型框架之外的考量因素，市场对汇率稳定诉求的定价尚不充分。
  \item 稀土磁材定价权正从上游氧化物向下游磁材环节转移，成本传导能力改善使磁材企业盈利模型发生结构性变化。
\end{itemize}
\section{宏观与利率}
\textbf{G10消费者健康度分化加剧，北美实际收入增长放缓是核心信号，欧元区相对稳健，日本表现最佳。}

\begin{itemize}
  \item 美国消费者信心跌至历史冰点，实际收入增长已降至历史最低5\%分位，能源价格上升正在侵蚀实际购买力。
  \item 美国消费韧性正从收入驱动转向储蓄消耗，储蓄率已降至2022年6月以来最低，预计未来几个月支出将放缓。
  \item 欧元区消费者健康度处于P65分位，劳动力市场仍紧张，资产负债表健康，是G10中最稳健的区域。
  \item 日本消费者健康度处于P50分位，基础工资增速仍超3\%，劳动力市场接近历史低位，表现最佳。
  \item 加拿大是G10中最脆弱的节点，消费者健康度仅P30分位，实际收入增长处于P20低位，房价持续下跌。
  \item 英国消费者健康度偏弱（P40分位），劳动力市场疲软，消费者信心低迷。
  \item 澳大利亚消费者健康度处于中位（P50），但名义消费环比下滑，运输支出受中东局势影响。
\end{itemize}
\begin{figure}[htbp]
\centering
\includegraphics[width=0.88\linewidth]{figures/fig_009.jpg}
\caption{Exhibit 1 - Exhibit 1: Balance Sheets Remain a}
\end{figure}
\begin{figure}[htbp]
\centering
\includegraphics[width=0.88\linewidth]{figures/fig_010.jpg}
\caption{Exhibit 2 - Exhibit 2: Real Income Growth Slowed in Canada but Improved in Japan; Real Spending Growth Worsened in Australia but Edged Up in Canada Change in Percentile}
\end{figure}
\begin{figure}[htbp]
\centering
\includegraphics[width=0.88\linewidth]{figures/fig_011.jpg}
\caption{Exhibit 3 - Exhibit 3: Our Consumer Health Indicator Stands Around P45 in the US and P25 in Canada Consumer Health Indicators: North America}
\end{figure}
\begin{figure}[htbp]
\centering
\includegraphics[width=0.88\linewidth]{figures/fig_012.jpg}
\caption{Exhibit 4 - Exhibit 4: Our Consumer Health Indicator Stands Around P65 in the Euro Area and P35 in the UK Consumer Health Indicators: Western Europe}
\end{figure}
\begin{figure}[htbp]
\centering
\includegraphics[width=0.88\linewidth]{figures/fig_013.jpg}
\caption{Exhibit 5 - Exhibit 5: Our Consumer Health Indicator Stands Around P50 in Japan and Australia Consumer Health Indicators: Asia Pacific}
\end{figure}
\begin{figure}[htbp]
\centering
\includegraphics[width=0.88\linewidth]{figures/fig_018.jpg}
\caption{Exhibit 10 - Exhibit 10: The Saving Rate Ticked Down in the US and Canada}
\end{figure}
\begin{figure}[htbp]
\centering
\includegraphics[width=0.88\linewidth]{figures/fig_019.jpg}
\caption{Exhibit 11 - Exhibit 11: The Unemployment Rate Ticked Up in the UK, Canada Australia but Ticked Down in Japan and the Euro Area}
\end{figure}
\begin{figure}[htbp]
\centering
\includegraphics[width=0.88\linewidth]{figures/fig_020.jpg}
\caption{Exhibit 12 - Exhibit 12: Wage Growth Accelerated in Japan but Ticked Down from Elevated Levels (Due to Compositional Distortions) in Canada}
\end{figure}
\begin{figure}[htbp]
\centering
\includegraphics[width=0.88\linewidth]{figures/fig_021.jpg}
\caption{Exhibit 13 - Exhibit 13: Morning Consult's High Frequency Measure of Consumer Sentiment Recovered Slightly in April but Remains Depressed Amid the War in the Middle East Consumer Sentiment Surveys}
\end{figure}
\begin{figure}[htbp]
\centering
\includegraphics[width=0.88\linewidth]{figures/fig_022.jpg}
\caption{Exhibit 14 - Exhibit 14: Sentiment Is Weakest in Australia for Both Low- and High-Income Consumers Consumer Sentiment by Income Group}
\end{figure}
\begin{figure}[htbp]
\centering
\includegraphics[width=0.88\linewidth]{figures/fig_023.jpg}
\caption{Exhibit 15 - Exhibit 15: Interest Expenses Remain Elevated in Canada and Australia}
\end{figure}
\begin{figure}[htbp]
\centering
\includegraphics[width=0.88\linewidth]{figures/fig_024.jpg}
\caption{Exhibit 16 - Exhibit 16: House Prices Continue to Decline in Canada}
\end{figure}
{\small\color{refgray}\textbf{References:} [R004] G10消费者健康正在分化，北美实际收入增长放缓才是核心信号}

\section{AI/算力与储能}
\textbf{全球储能产业正经历供给端结构性重组，AI数据中心成为储能新引擎，但真正改变行业格局的是产能重新布局与技术路线切换。}

\begin{itemize}
  \item 美国AI数据中心正在成为储能新引擎，LG Energy Solution签下16亿美元大单，为DTE Energy供应6GWh储能电池，专门支持AI数据中心相关电网项目。
  \item 全球储能产业正从规模竞赛进入成本与场景的双重分化期，企业通过轻资产化策略（如LGES出售合资工厂资产）换取现金流与灵活性。
  \item 中国LFP材料与换电模式双线推进，天赐材料投资21亿元建LFP正极材料工厂，宁德时代在粤港澳大湾区启动轻型电卡换电网络。
  \item 换电模式从乘用车延伸到商用车，BaaS（电池即服务）在重卡领域开始落地，运营成本比柴油车低约50\%。
  \item 电池行业当前周期与2021-2022年超级周期本质不同，价格弹性远小于上一轮，但盈利改善逻辑更健康，市场焦点应从价格修复转向出货量驱动的现金流增长。
\end{itemize}
\begin{figure}[htbp]
\centering
\includegraphics[width=0.88\linewidth]{figures/fig_001.jpg}
\caption{Figure 1 - \textbackslash{}- Range-bound trading in 2H26-1H27. This argues for a more range-bound trading approach for material names and tier-2 battery makers. In 2025, LFP cathode names such as Hunan Yuneng worked well as price-recovery trades, followed by separator names such as Yun}
\end{figure}
{\small\color{refgray}\textbf{References:} [R001] 全球储能正在经历一场“供给重构”，而非简单的需求回暖; [R002] 电池行业的两个周期：市场真正低估的不是需求弹性，而是盈利兑现的可持续性}

\section{能源与大宗}
\textbf{铜价核心支撑来自非美市场供给收紧与美国进口激增导致的库存地理错配，稀土磁材定价权正从上游向下游转移。}

\begin{itemize}
  \item 铜价当前及未来两年的核心支撑来自非美市场供给持续低于预期，以及美国进口激增导致的全球铜库存地理错配。
  \item 2026年全球矿山供给预测被削减约35万吨，主要因印尼Grasberg和刚果金Kamoa-Kakula矿山复产慢于预期，预计要到2028年才能恢复满产。
  \item 非美市场铜供需缺口从此前6万吨/4万吨大幅上调至64万吨/17万吨，美国在囤铜，其他市场承受更紧的供应。
  \item 中国废铜产出同比下降12\%，因增值税发票合规趋严，废铜替代精炼铜的能力受限，精炼铜供需弹性更低。
  \item 稀土磁材定价权正从上游稀土氧化物向下游磁材环节转移，金力永磁通过成本加成模式和客户锁价合同有效传导上游涨价压力。
  \item 2026年一季度氧化镨钕价格同比上涨约50\%，而金力永磁毛利率反向扩张约6个百分点至22\%，显示磁材行业盈利模型发生系统性改善。
  \item 稀土磁材需求增长亮点来自工业机器人（销售额同比增超80\%）和电动重卡，电动车仍是主力但增速趋稳。
\end{itemize}
\begin{figure}[htbp]
\centering
\includegraphics[width=0.88\linewidth]{figures/fig_002.jpg}
\caption{Exhibit 1 - Exhibit 1: Two-Sided Risks to our LME Copper Price Forecast LME Copper Price}
\end{figure}
\begin{figure}[htbp]
\centering
\includegraphics[width=0.88\linewidth]{figures/fig_003.jpg}
\caption{Exhibit 2 - Exhibit 2: We Expect the Ex-US Market to be in Substantial Deficit This Year, Keeping the LME Price Supported Global Copper Market Balance}
\end{figure}
\begin{figure}[htbp]
\centering
\includegraphics[width=0.88\linewidth]{figures/fig_004.jpg}
\caption{Exhibit 3 - Exhibit 3: We Have Removed \textbackslash{}\textasciitilde{}350kt of Mine Supply in 2026 As the Impacts of Last Year's Mine Supply Disruptions Persist Change in GS Copper Mine Supply Forecast (May vs. February 2026)}
\end{figure}
\begin{figure}[htbp]
\centering
\includegraphics[width=0.88\linewidth]{figures/fig_005.jpg}
\caption{Exhibit 4 - Exhibit 4: China Copper Scrap Output Falls 12\% YoY YTD on Invoice Scrutiny China Domestic Copper Scrap Output}
\end{figure}
\begin{figure}[htbp]
\centering
\includegraphics[width=0.88\linewidth]{figures/fig_006.jpg}
\caption{Exhibit 5 - Exhibit 5: We Expect US Copper Inventory Growth to Accelerate in the Coming Weeks, but Slow in H2 Under the Base Case of No Tariff Announcement US Copper Inventory Change}
\end{figure}
\begin{figure}[htbp]
\centering
\includegraphics[width=0.88\linewidth]{figures/fig_007.jpg}
\caption{Exhibit 6 - Exhibit 6: We Estimate That a 1-Day Decline in the Copper Balance (Inventory Days) Results in a \textbackslash{}\textasciitilde{}1.4\% Boost to the LME Price}
\end{figure}
\begin{figure}[htbp]
\centering
\includegraphics[width=0.88\linewidth]{figures/fig_008.jpg}
\caption{Exhibit 7 - Exhibit 7: Chinese Apparent Copper Demand Remained Resilient at \$+1\textbackslash{}\%\$ YoY in Q1 2026 China Copper Apparent Demand}
\end{figure}
\begin{figure}[htbp]
\centering
\includegraphics[width=0.88\linewidth]{figures/fig_039.jpg}
\caption{Exhibit 1 - Exhibit 1: Total 1Q26 revenue increased by \textbackslash{}\textasciitilde{}16\% YoY to \textbackslash{}\textasciitilde{}RMB2bn}
\end{figure}
\begin{figure}[htbp]
\centering
\includegraphics[width=0.88\linewidth]{figures/fig_040.jpg}
\caption{Exhibit 2 - Exhibit 2: Overseas revenue increased \$\textbackslash{}sim 22\textbackslash{}\%\$ YoY}
\end{figure}
\begin{figure}[htbp]
\centering
\includegraphics[width=0.88\linewidth]{figures/fig_041.jpg}
\caption{Exhibit 3 - Exhibit 3: Gross margins expanded by \$\textbackslash{}sim 6\textbackslash{}\%\$ to \$\textbackslash{}sim 22\textbackslash{}\%\$}
\end{figure}
\begin{figure}[htbp]
\centering
\includegraphics[width=0.88\linewidth]{figures/fig_042.jpg}
\caption{Exhibit 4 - Exhibit 4: Chinese NdFeB Exports increased \textbackslash{}\textasciitilde{}5\% YoY in 1Q26 NdFeB products, annualised (ktpa)}
\end{figure}
\begin{figure}[htbp]
\centering
\includegraphics[width=0.88\linewidth]{figures/fig_043.jpg}
\caption{Exhibit 6 - Exhibit 6: We forecast NdPrO deficits over the medium term to 2027 and perhaps beyond... GS Rare Earth (NdPrO) SD model - net market balance NdPrO (kt)}
\end{figure}
\begin{figure}[htbp]
\centering
\includegraphics[width=0.88\linewidth]{figures/fig_044.jpg}
\caption{Exhibit 5 - Exhibit 5: China's NdFeB alloy magnet export prices increased in 1Q26, reflecting higher RE oxide prices. NdFeB Magnet Export Prices (US\textbackslash{}\$/kg)}
\end{figure}
\begin{figure}[htbp]
\centering
\includegraphics[width=0.88\linewidth]{figures/fig_045.jpg}
\caption{Exhibit 7 - Exhibit 7: Growth in Electric Vehicles (EV), robotics, variable frequency air conditioning (VFAC), and mobile phones/consumer electronics should nearly double NdPrO demand by 2035 NdPrO demand by}
\end{figure}
\begin{figure}[htbp]
\centering
\includegraphics[width=0.88\linewidth]{figures/fig_046.jpg}
\caption{Exhibit 8 - Exhibit 8: LYC, MP Materials, ILU and emerging companies to drive supply growth over the next 5-years NdPrO supply by}
\end{figure}
\begin{figure}[htbp]
\centering
\includegraphics[width=0.88\linewidth]{figures/fig_047.jpg}
\caption{Exhibit 8 - Exhibit 10: GS Rare Earth (NdPrO) SD model}
\end{figure}
{\small\color{refgray}\textbf{References:} [R003] 铜价真正的支撑不在需求故事，而在“非美市场”的供给收紧与库存错配; [R007] 稀土磁材的定价权正在从上游向中游转移}

\section{区域市场与地缘政治}
\textbf{中东局势缓和情景下，南非、希腊、波兰等具备高油价负相关、低利率敏感与超卖三重特征的市场反弹弹性最大。}

\begin{itemize}
  \item 中东局势缓和交易的核心逻辑不在能源，而在高油价负相关、低本地利率敏感与深度超卖三重特征的叠加效应。
  \item 南非是最大受益者，矿业（铂金、黄金）、金融（银行、保险）和消费（零售、日用品）板块反弹弹性最强。
  \item 希腊和波兰在缓和情景下也将显著受益，土耳其银行也在看多名单中。
  \item 沙特和东欧能源、化工、公用事业板块短期压力较大，但沙特中长期仍受油价高位支撑。
  \item 人民币中间价模型显示政策当局正系统性引入模型框架之外的考量因素，市场对汇率稳定诉求的定价尚不充分。
  \item 模型误差从2025年初的-1800个基点收窄至2026年4月的+600个基点，央行操作越来越可预测。
\end{itemize}
\begin{figure}[htbp]
\centering
\includegraphics[width=0.88\linewidth]{figures/fig_034.jpg}
\caption{Exhibit 1 - Exhibit 1: What if the Strait reopens? From a technical and tactical perspective, our analysis suggests this would be most positive for S. Africa, Greece, Poland and the UAE MSCI EEMEA: Combined Tactical Score in De-escalation Scen}
\end{figure}
\begin{figure}[htbp]
\centering
\includegraphics[width=0.88\linewidth]{figures/fig_035.jpg}
\caption{Exhibit 2 - Exhibit 2: EEMEA equities are attempting to break out of their recent downtrend resistance level MSCI EEMEA: Index Price (USD)}
\end{figure}
\begin{figure}[htbp]
\centering
\includegraphics[width=0.88\linewidth]{figures/fig_036.jpg}
\caption{Exhibit 4 - Exhibit 4: Bearish: Applying the same filters as above and looking at the opposite end yields a list with a high concentration of Energy, Chemicals, Utilities sectors across Saudi Arabia and EM Europe Exhibit 5: Absolute performanc}
\end{figure}
\begin{figure}[htbp]
\centering
\includegraphics[width=0.88\linewidth]{figures/fig_037.jpg}
\caption{Exhibit 7 - Exhibit 7: EEMEA equities are currently extremely oversold vs EM ...}
\end{figure}
\begin{figure}[htbp]
\centering
\includegraphics[width=0.88\linewidth]{figures/fig_038.jpg}
\caption{Exhibit 6 - Exhibit 6: ... And relative to MSCI EM}
\end{figure}
{\small\color{refgray}\textbf{References:} [R005] 中东局势缓和：市场真正低估的不是油价，而是南非的反弹弹性; [R006] 人民币中间价模型正在发出一个被忽视的信号}

\section*{结语}
综合来看，当前市场主线正从需求侧叙事转向供给侧结构性分化，储能、铜、稀土等板块的定价逻辑正在重塑，而G10消费者健康度的分化与地缘政治缓和情景下的区域市场机会，为投资者提供了多维度的观察视角。
\vfill
{\small\color{refgray}Personal reading notes and learning share only. Not investment advice.}
\end{document}
